% ============================================================
% Chapter 5: Discussion
% ============================================================

\chapter{Discussion}
\label{chap:discussion}

This chapter interprets the findings from Chapter~\ref{chap:results} in relation to prior work, discusses the pricing and access implications of the geographic gradient, and examines the threats to validity.

% ------------------------------------------------------------
\section{Interpretation of RQ1: Scope of Adoption}
\label{sec:disc:rq1}
% ------------------------------------------------------------

The 74.4\% adoption rate is both high in absolute terms and closely consistent with the most comparable prior work. Robbes et al.~\cite{robbesAgenticVery2026} reported 71--76\% adoption in newly created repositories across all GitHub languages. The JS/TS/npm corpus studied here falls at the upper end of that range. Two factors plausibly explain why the JS/TS/npm ecosystem sits at the high end rather than the low end of the previously documented range.

First, the corpus was deliberately restricted to the npm ecosystem, which is the most tooling-intensive slice of JavaScript development. npm projects maintain explicit lockfiles, versioned dependency lists, and scripts in \texttt{package.json} --- all practices associated with engineering discipline and openness to new tooling. The ecosystem-specificity of this study means the corpus population is predisposed to tool adoption relative to a random sample of all GitHub JS/TS repositories.

Second, Claude Code's design aligns directly with the npm workflow. It is a terminal-first, CLI-native agent that integrates naturally into an \texttt{npm run} environment. This structural affinity between the most-adopted agent (Claude Code, 63.7\% of adopted repositories) and the ecosystem under study likely inflates the JS/TS/npm rate relative to a cross-language average.

The 50.5\% adoption rate in February 2025 rising to 84.3\% by March 2026 traces a rapid growth trajectory compressed into 13 months. The acceleration between December 2025 and February 2026 (+11 percentage points in two months) coincides with the period of intense competition between Claude Code, OpenAI Codex CLI, and Gemini CLI, suggesting that multiple competing tools reaching maturity simultaneously accelerates aggregate adoption. The plateau at approximately 82--84\% from February to May 2026 is consistent with a natural saturation ceiling: the fraction of developers who actively create public npm projects and are unwilling or unable to use AI agents appears to have stabilised near 16--18\%.

% ------------------------------------------------------------
\section{Interpretation of RQ2: Signal Undercounting}
\label{sec:disc:rq2}
% ------------------------------------------------------------

The 27.21 percentage-point undercounting gap between single-signal detection (47.2\%) and the multi-signal pipeline (74.4\%) is the central methodological finding of this study. It is larger than the 6.46~pp gap reported by Robbes et al.~\cite{robbesAgenticMuch2026}, which compared two heuristic sets within a single-pass framework. The difference is explained by the structure of the comparison: Robbes et al.\ compared two overlapping heuristic variants, whereas this study compares a single signal family (configuration files) against four qualitatively distinct passes that each recover a disjoint population of adopted repositories.

The 16,364 repositories that the full pipeline detects but single-signal detection misses divide into at least two distinct populations. The largest group (11,937 repositories from Pass~2) are developers who actively use agents via commit trailers but never created a configuration file at the repository root --- either because they use agents interactively without project-level setup or because they use agents that do not follow the config-file convention. The second group (70 repositories from Pass~3 plus those upgraded by Stage~3d) are captured only through deep signals such as PR body attribution, bot-authored commits, or npm SDK dependencies. Both populations are systematically invisible to any study that relies solely on configuration files.

The McNemar $c$ value of 40 --- the number of repositories that single-signal catches but multi-signal misses --- is small but non-zero. These 40 repositories had a configuration file that was detected in Pass~1 but whose label was later revised or whose repository was excluded from the final corpus on data-quality grounds. The asymmetry between $b = 16{,}404$ and $c = 40$ is so extreme (ratio 410:1) that it confirms the multi-signal pipeline is essentially a strict superset of single-signal detection.

The finding aligns with Khosravani and Mockus~\cite{khosravani2026}, who documented a potentially 30-fold undercounting factor for single-signal approaches in a different language context. The gap documented here (a factor of approximately 1.58$\times$ in raw count, $74.4\% / 47.2\%$) is well within the range that prior work suggested, and the direction of the bias is unambiguous: single-signal approaches undercount and never overcount adoption.

% ------------------------------------------------------------
\section{Interpretation of RQ3: Geographic Variation}
\label{sec:disc:rq3}
% ------------------------------------------------------------

The 63.5-percentage-point spread between the highest-adopting country (Saudi Arabia, 90.9\%) and the lowest (Rwanda, 27.4\%) documents a geographic adoption divide that is both statistically robust and substantively large. The chi-squared test rejects homogeneity with $\chi^2 = 2{,}577.03$ on 65~degrees of freedom, and 45 of the 66 eligible countries are individually significant against the global mean.

The gradient broadly follows the World Bank income classification: all 14 countries in the significantly high-adoption group are high-income or upper-middle-income, and all 20 countries in the significantly low-adoption group are lower-middle-income or low-income. This pattern is consistent with the income-stratified adoption gradient documented by Daniotti et al.~\cite{daniotti2025} for Python and with the GDP--contribution relationship documented by Wachs et al.~\cite{wachs2022} for open-source participation more broadly.

Several countries deviate from income-based predictions in ways that are worth noting. South Korea (rank~4, 86.9\%) and Taiwan (rank~5, 85.0\%) both adopt at rates higher than their income classification alone would predict: both countries have concentrated, engineering-intensive technology sectors with high English-language proficiency and strong integration into the global JavaScript ecosystem. The Gulf states (Saudi Arabia rank~1, UAE rank~3) show similarly elevated rates despite having smaller developer communities on GitHub --- suggesting that in those countries, the developers who are on GitHub are particularly likely to be early adopters of AI tooling. By contrast, China (rank~30, 76.9\%) adopts below the rate that its upper-middle-income classification would predict when compared against countries like South Korea or Malaysia, possibly reflecting a combination of VPN overhead, partial access constraints, and the moderate share of domestic tools in the Chinese agentic ecosystem (2.3\% of Chinese agentic repositories use a Chinese domestic tool).

India's position (rank~54, 52.9\%) merits attention. With 4,710 repositories, India is the third-largest country in the located corpus, and its below-average adoption rate exerts a large downward pull on the South Asian regional figure. The gap between the United States (82.9\%) and India (52.9\%) --- 30 percentage points --- within a single ecosystem represents one of the most consequential geographic disparities in the data.

% ------------------------------------------------------------
\section{Pricing, Access, and the Geographic Gradient}
\label{sec:disc:pricing}
% ------------------------------------------------------------

The income-group pricing analysis provides the most direct evidence that pricing is a mediating factor in the geographic adoption gradient. As income group falls, the share of adopted repositories using free tools rises from 11.4\% in high-income countries to 24.3\% in lower-middle-income countries. The low-income group (27.5\% adoption) uses free tools in 41.7\% of its adopted repositories.

This pattern is consistent with the finding of Daniotti et al.~\cite{daniotti2025} that adoption gaps between high-income and lower-income developer communities may widen over time if paid tools continue to outpace free alternatives in capability. The JS/TS/npm data supports this concern: at the tool level, Claude Code (the most capable paid agent) accounts for 63.7\% of adopted repositories globally but a smaller fraction in lower-income regions. The free-tool dominance in LMIC and LIC countries is not driven by conscious ecosystem preference but by economic constraint.

There are reasons, however, to be cautious about over-interpreting the pricing gradient. Free tools such as GitHub Copilot (free tier), Gemini CLI, and OpenCode account for a significant fraction of adoption even in high-income countries, so the high paid-tool share in HIC is partly a supply-side phenomenon: Claude Code and Cursor are simply more widely used there regardless of price sensitivity. Disentangling pricing effects from capability preferences requires a study design beyond what observational repository data can support.

% ------------------------------------------------------------
\section{Threats to Validity}
\label{sec:disc:threats}
% ------------------------------------------------------------

\subsection{Internal Validity}

The most significant internal threat is the \textbf{arbitrariness of labelling thresholds}. The boundary between \emph{agentic} and \emph{AI-assisted} is a design choice: requiring one STRONG signal for agentic classification is defensible but not uniquely correct. A stricter threshold (requiring two STRONG signals) would lower the agentic rate; a looser threshold (treating any MEDIUM signal as sufficient for agentic) would raise it. The adoption rate (74.4\%) is robust to this choice because both classes contribute to it; the split between agentic (67.1\%) and AI-assisted (7.3\%) is not. Future work should test sensitivity to threshold variation.

The \textbf{geographic inference pipeline} introduces a second internal threat. With 43.4\% of owners unattributed, the located subset may be systematically biased. If the propensity to leave location information blank is higher in certain countries, those countries will be under-represented in the geographic analysis. The pipeline mitigates this bias through four supplementary passes, but it cannot eliminate it entirely. The fact that China required the most aggressive supplementary resolution (the deep pass added 1,941 Chinese owners, a 127\% increase) illustrates how the bias varies by country.

\subsection{External Validity}

The corpus covers only public GitHub repositories. Private repositories --- which may account for a substantial fraction of commercial development --- are entirely unobservable. If AI agent adoption rates differ between public and private repositories (plausibly in either direction: public repositories may adopt more to signal modernity, or private ones may adopt more because commercial resources are available), the 74.4\% rate cannot be generalised to all JS/TS/npm development.

The npm ecosystem is the world's largest package registry but is not coextensive with all JavaScript or TypeScript development. Projects using server-side TypeScript without npm (some Deno projects, Cloudflare Workers with custom runtimes, enterprise Java--TypeScript hybrid stacks) are outside the corpus. The vanilla category partially captures non-npm JS/TS projects, but it is itself filtered by the same star and commit criteria.

\subsection{Construct Validity}

Detecting agent signals does not establish that an agent was used productively or that it contributed to the project's quality. A CLAUDE.md file may be present because a developer copied a template; a \texttt{Co-authored-by: Claude} trailer may appear in one commit out of a thousand. The adoption label captures the \emph{presence} of agent-related artefacts, not the \emph{quality} or \emph{intensity} of use. The commit density analysis (Section~\ref{sec:res:extended}) partially addresses this by showing that 27.0\% of trailer-detected repositories are classified as pervasive users ($>$20\% of commits), but this does not resolve the ambiguity for the larger group of repositories classified via configuration files.

% ------------------------------------------------------------
\section{Implications}
\label{sec:disc:implications}
% ------------------------------------------------------------

\textbf{For tool vendors.}
The terminal-first, CLI-native design of Claude Code appears to be a structural advantage in the npm ecosystem, where \texttt{npm run} and \texttt{npx} workflows are native. The 63.7\% penetration among adopted repositories suggests that the design decision to build an agent that integrates with the shell rather than the IDE has paid off in the ecosystem where the shell is a primary interface. Competing tools targeting JS/TS development would benefit from similar CLI-native integration.

\textbf{For educators and developers.}
The adoption timing data (55.9\% immediate, median 3.0 days) suggests that AI agent adoption is now a project-initiation behaviour rather than a workflow addition to an existing project. Developer education programmes that introduce AI agents as part of project bootstrapping --- rather than as an add-on to established workflows --- align with the actual adoption pattern.

\textbf{For open-source foundations and policy.}
The pricing gradient documents a measurable access gap: developers in lower-income countries adopt AI agents at half the rate of high-income developers, and when they do adopt, they disproportionately use free tools whose capabilities may be more limited. Open-source foundations and AI labs with geographic equity programmes should consider this data when designing free-tier offerings, regional pricing, or subsidy schemes. The Gulf--Iran divide within MENA (87.9\% vs.\ 47.1\%) shows that even within a single region, national-level differences in economic access and internet infrastructure translate directly into adoption gaps.
